\appendix

\section{Proof Details}
\label{sec:proof-details}

\subsection{Legality of the Canonical Rounded Execution}
\label{app:step-001}

\begin{lemma}[Pathwise legality of the canonical rounded policy]
\label{lem:step-001-canonical-validity}
Under Assumptions~\ref{assump:sq-parameter-regime} and
\ref{assump:universal-adversarial-sq}, fix a distribution \(\mathcal D\), a
target \(h\in\mathcal H\), and a tape \(r\in\Omega\).  Generate the
interaction recursively by replying to each actually reached bounded query
\(q_t\) with
\[
v_t:=\mathbb E_{x\sim\mathcal D}q_t(x,h(x)),
\qquad
a_t:=\rho(v_t).
\]
Then every reached reply satisfies
\[
|a_t-v_t|\le\frac{1}{K}\le\tau,
\qquad
K=\left\lceil\frac{1}{\tau}\right\rceil.
\]
Consequently, the recursion is an actual tolerance-valid execution of some
length \(T\le m\).  This includes early stopping and the empty execution when
\(m=0\), for every finite \(\tau>0\), every fixed midpoint rule, and every
tape \(r\), without a condition on the tape law.
\end{lemma}

\begin{proof}
Because \(\tau>0\), the integer \(K=\lceil1/\tau\rceil\) is finite and
positive.  The grid in \eqref{eq:rounding-grid} has adjacent spacing \(2/K\)
and covers \([-1,1]\).  A point between two consecutive grid points is
within half that spacing of at least one endpoint.  Thus, for every
\(v\in[-1,1]\),
\begin{equation}
|\rho(v)-v|\le\frac{1}{K}.
\label{eq:step-001-grid-cover}
\end{equation}
At a midpoint both adjacent points are at distance \(1/K\), so the fixed tie
rule preserves the bound.  Moreover,
\begin{equation}
K=\left\lceil\frac{1}{\tau}\right\rceil
\ge\frac{1}{\tau}>0
\quad\Longrightarrow\quad
\frac{1}{K}\le\tau.
\label{eq:step-001-radius-tolerance}
\end{equation}
This also covers \(\tau\ge1\): then \(K=1\), the grid is
\(\{-1,1\}\), and its covering radius is \(1\le\tau\).

Fix the chosen \((\mathcal D,h,r)\).  Once the tape is fixed, the learner's
next action is deterministic given the replies actually received.  Begin
with the empty reached history.  If the learner stops, the interaction is
complete.  Otherwise, it issues a first query
\(q_1:\mathcal X\times\{+1,-1\}\to[-1,1]\).  Pointwise boundedness gives
\[
-1\le q_1(x,h(x))\le1,
\]
and hence \(v_1\in[-1,1]\).  Equations
\eqref{eq:step-001-grid-cover}--\eqref{eq:step-001-radius-tolerance} give
\( |a_1-v_1|\le\tau\), exactly the additive condition
\eqref{eq:sq-validity}.

For the induction step, suppose the first \(t-1\) canonical replies form the
reached valid prefix.  That prefix and the fixed tape determine whether the
learner stops or issues a next bounded query \(q_t\).  In the latter case,
the same pointwise bound gives \(v_t\in[-1,1]\), regardless of how \(q_t\)
depends on the tape and earlier replies.  Therefore
\begin{equation}
|a_t-v_t|=|\rho(v_t)-v_t|
\le\frac{1}{K}\le\tau,
\label{eq:step-001-reached-validity}
\end{equation}
and appending \(a_t\) yields the next valid reached prefix.  The induction
continues only until the learner stops, and
Assumption~\ref{assump:universal-adversarial-sq} bounds the reached length by
\(T\le m\).  When \(T\ge1\), the pathwise conclusion may equivalently be
written
\[
\max_{1\le t\le T}|a_t-v_t|\le\tau.
\]
No bound on \(\sum_t|a_t-v_t|\) is required or asserted.  Adaptive changes
of later queries therefore create no accumulated response defect.

If the learner stops after a valid prefix, no unreached query requires a
reply.  If \(m=0\), no query can be reached, so the empty interaction is the
actual tolerance-valid execution by vacuity.  The argument is pointwise in
\(r\) and uses neither \(\mu\), atoms, nor finite support of the tape space.
The oracle uses \(v_t\) only to compute \(a_t\); the learner receives
\(a_t\), not \(v_t\).  Finally, the recursion follows only the history
generated by these canonical replies.  It does not replay a prescribed
tolerance-invalid string or assign such a string a terminal output.
\end{proof}

\begin{proposition}[Exact membership on the actual canonical execution]
\label{prop:step-001-canonical-membership}
Under Assumptions~\ref{assump:sq-parameter-regime},
\ref{assump:universal-adversarial-sq}, and
\ref{assump:canonical-rounded-output-catalog}, together with
Lemma~\ref{lem:step-001-canonical-validity}, for every distribution
\(\mathcal D\), target \(h\in\mathcal H\), and tape \(r\in\Omega\),
\[
A_r^{\mathcal O^{\rho}_{\mathcal D,h}}(\mathcal D,h)
\in\mathcal C_A^{\rho}=(g_1,\ldots,g_L)
\]
as equality of functions on \(\mathcal X\).  The catalog is the same ordered
catalog fixed before the instance.  The statement includes the no-query
execution and makes no claim about another valid policy or a prescribed
tolerance-invalid response string.
\end{proposition}

\begin{proof}
Fix \((\mathcal D,h,r)\).  By
Lemma~\ref{lem:step-001-canonical-validity}, the recursively generated
interaction against \(\mathcal O^{\rho}_{\mathcal D,h}\) is an actual
tolerance-valid canonical execution, including early stopping and the empty
execution.  The antecedent of
Assumption~\ref{assump:canonical-rounded-output-catalog} is therefore met.
That assumption supplies an index \(i\in\{1,\ldots,L\}\), not necessarily
unique when catalog entries repeat, such that
\begin{equation}
A_r^{\mathcal O^{\rho}_{\mathcal D,h}}(\mathcal D,h)(x)=g_i(x)
\qquad\text{for every }x\in\mathcal X.
\label{eq:step-001-function-equality}
\end{equation}
This is zero-residual equality of terminal functions, rather than approximate
membership.  Since the fixed triple was arbitrary, the conclusion holds for
all instances and tapes.  No measurable selection of \(i\) as a function of
\(r\) is made, so the argument is unchanged for an arbitrary or nonatomic
tape space.

Together with Lemma~\ref{lem:step-001-canonical-validity}, this proves that
every canonical interaction used below is both an actual valid execution and
catalog-valued.  The scope remains canonical-policy-only: no conclusion is
drawn for another valid policy, a synthetic invalid transcript, or an
all-policy output range.
\end{proof}

\begin{proof}[Conclusion of the canonical-execution argument]
Let \((\mathcal D,h)\) be arbitrary.  Lemma
\ref{lem:step-001-canonical-validity} applies pathwise to every tape: each
reached response satisfies \(|a_t-v_t|\le1/K\le\tau\), so the adaptive
recursion is an actual valid execution of length at most \(m\).  This
includes early stopping, the actual no-query execution at \(m=0\), fixed
midpoint ties, every \(\tau>0\), and an arbitrary tape space.  Proposition
\ref{prop:step-001-canonical-membership} then gives
\[
\forall r\in\Omega,
\qquad
A_r^{\mathcal O^{\rho}_{\mathcal D,h}}(\mathcal D,h)
\in\mathcal C_A^{\rho}
\]
as equality of functions on \(\mathcal X\).  Since \((\mathcal D,h)\) was
arbitrary, these conclusions provide exactly the actual canonical executions
and fixed-catalog membership used later, without defining an output on an
invalid string or extending the catalog to another policy.
\end{proof}

\subsection{Selector-Free Extraction of a Correlated Catalog Function}
\label{app:step-002}

\begin{lemma}[Attainment of a low-risk actual catalog output]
\label{lem:step-002-attained-low-risk}
Under Assumptions~\ref{assump:sq-parameter-regime},
\ref{assump:universal-adversarial-sq}, and
\ref{assump:canonical-rounded-output-catalog}, together with
Lemma~\ref{lem:step-001-canonical-validity} and
Proposition~\ref{prop:step-001-canonical-membership}, fix a distribution
\(\mathcal D\) and target \(h\in\mathcal H\).  For each \(r\in\Omega\), set
\[
G_r:=A_r^{\mathcal O^{\rho}_{\mathcal D,h}}(\mathcal D,h),
\qquad
Z(r):=\mathcal L_{\mathcal D,h}(G_r).
\]
Then \(Z:\Omega\to[0,1]\) is \(\mu\)-measurable, its range is nonempty and
finite, and there are \(r_*\in\Omega\) and \(i_*\in\{1,\ldots,L\}\) such
that
\[
G_{r_*}=g_{i_*}\quad\text{on }\mathcal X,
\qquad
\mathcal L_{\mathcal D,h}(g_{i_*})=Z(r_*)\le\varepsilon.
\]
Occurrence means existence of the tape \(r_*\) on the actual canonical
execution; no atom or positive-mass output fiber is required.
\end{lemma}

\begin{proof}
Lemma~\ref{lem:step-001-canonical-validity} proves that the canonical policy
is tolerance-valid along the actual execution for every tape.  Instantiating
Assumption~\ref{assump:universal-adversarial-sq} at this one proved-valid
policy and the fixed \((\mathcal D,h)\) gives
\begin{equation}
Z:\Omega\to[0,1]\ \text{measurable},
\qquad
\mathbb E_{R\sim\mu}Z(R)\le\varepsilon.
\label{eq:step-002-scalar-risk}
\end{equation}
No random function, catalog coordinate, or output index is declared
measurable in obtaining \eqref{eq:step-002-scalar-risk}.

Proposition~\ref{prop:step-001-canonical-membership} gives, separately for
each tape,
\[
G_r\in\{g_1,\ldots,g_L\}
\]
as exact equality of functions.  Consequently, the actual terminal function
has at most \(L\) distinct values, and the scalar range
\[
S:=Z(\Omega)
\]
has at most \(L\) distinct real values.  This is a set-theoretic statement,
not a pushforward probability law or measurability claim for an index map.
Because \(\mu(\Omega)=1\), the set \(\Omega\), and hence \(S\), is nonempty.
Thus \(S\subseteq[0,1]\) is nonempty and finite.

Let \(s_*:=\min S\).  By the definition of the range, some
\(r_*\in\Omega\) satisfies \(Z(r_*)=s_*\).  Pointwise, \(Z\ge s_*\), so
positivity of the integral applied to \(Z-s_*\) gives
\begin{equation}
s_*=\int_{\Omega}s_*\,d\mu
\le\int_{\Omega}Z\,d\mu
=\mathbb E_{R\sim\mu}Z(R)
\le\varepsilon.
\label{eq:step-002-minimum-expectation}
\end{equation}
Only the bounded measurable scalar \(Z\) is integrated; no output function,
index, fiber indicator, or catalog-coordinate loss is integrated.

Apply Proposition~\ref{prop:step-001-canonical-membership} once at \(r_*\).
It gives an \(i_*\in\{1,\ldots,L\}\) such that \(G_{r_*}=g_{i_*}\)
pointwise.  Hence
\begin{equation}
\mathcal L_{\mathcal D,h}(g_{i_*})
=\mathcal L_{\mathcal D,h}(G_{r_*})
=Z(r_*)=s_*\le\varepsilon.
\label{eq:step-002-attained-risk}
\end{equation}
The choice is a single existential choice after
\((\mathcal D,h,r_*)\) is fixed.  Duplicate entries need not have a unique
index, and no selector over tapes, distributions, or targets is constructed.

The proof neither uses an atom of \(\mu\), finite tape support, nor positive
mass or measurability of an output fiber.  It does not require the singleton
\(\{r_*\}\) to be measurable.  For \(L=1\), the scalar range is a
singleton; duplicates only reduce its cardinality.  For \(m=0\), the actual
empty execution supplied by Lemma~\ref{lem:step-001-canonical-validity}
precedes \eqref{eq:step-002-scalar-risk}.  Every \(\tau>0\) remains covered.
If \(\varepsilon=0\), then \(0\le s_*\le0\), so the occurring catalog
function has risk zero.
\end{proof}

\begin{proposition}[Exact same-function conversion from risk to correlation]
\label{prop:step-002-exact-correlation}
Under Assumptions~\ref{assump:sq-parameter-regime},
\ref{assump:universal-adversarial-sq}, and
\ref{assump:canonical-rounded-output-catalog}, together with
Lemma~\ref{lem:step-002-attained-low-risk}, for every distribution
\(\mathcal D\) and target \(h\in\mathcal H\), some actually occurring
catalog function \(g_i\) satisfies
\begin{equation}
\mathbb E_{x\sim\mathcal D}[h(x)g_i(x)]
=1-2\mathcal L_{\mathcal D,h}(g_i)
\ge1-2\varepsilon
=\rho_{\varepsilon}.
\label{eq:step-002-correlation}
\end{equation}
The risk and correlation concern the same function, target, and distribution.
\end{proposition}

\begin{proof}
Fix \((\mathcal D,h)\), and let \(g_{i_*}\) be the function supplied by
Lemma~\ref{lem:step-002-attained-low-risk}.  Since
\(h(x),g_{i_*}(x)\in\{+1,-1\}\), their product is \(+1\) when they agree and
\(-1\) when they disagree.  Pointwise,
\begin{equation}
h(x)g_{i_*}(x)
=1-2\mathbf 1\{g_{i_*}(x)h(x)<0\}.
\label{eq:step-002-binary-identity}
\end{equation}
The indicator is exactly the event in \eqref{eq:binary-risk}.  Taking the
\(\mathcal D\)-expectation in \eqref{eq:step-002-binary-identity} gives
\begin{equation}
\mathbb E_{x\sim\mathcal D}[h(x)g_{i_*}(x)]
=1-2\Pr_{x\sim\mathcal D}[g_{i_*}(x)h(x)<0]
=1-2\mathcal L_{\mathcal D,h}(g_{i_*}).
\label{eq:step-002-risk-correlation-identity}
\end{equation}
By \eqref{eq:step-002-attained-risk}, the risk is at most \(\varepsilon\).
Multiplication by \(-2\) reverses the inequality, so
\[
1-2\mathcal L_{\mathcal D,h}(g_{i_*})
\ge1-2\varepsilon=\rho_{\varepsilon}.
\]
There is no approximation, rounding term, change of function, or
distributional transfer.  At \(\varepsilon=0\), the same argument gives
correlation exactly one.  Since the fixed pair was arbitrary, this proves the
stated per-distribution catalog response without a measurable selector,
output law, positive-mass fiber, or output from another policy.
\end{proof}

\begin{proof}[Conclusion of the scalar-loss extraction argument]
Fix an arbitrary \((\mathcal D,h)\).  Lemma
\ref{lem:step-001-canonical-validity} makes the canonical policy valid, and
Proposition~\ref{prop:step-001-canonical-membership} makes its terminal range
finite.  Lemma~\ref{lem:step-002-attained-low-risk} then uses only the
measurable scalar actual-run loss to produce an actually occurring
\(g_{i_*}\) with
\(\mathcal L_{\mathcal D,h}(g_{i_*})\le\varepsilon\).
Proposition~\ref{prop:step-002-exact-correlation} applies the pointwise binary
identity to that same function and gives
\[
\mathbb E_{x\sim\mathcal D}[h(x)g_{i_*}(x)]
\ge1-2\varepsilon=\rho_{\varepsilon}.
\]
Arbitrariness of the pair yields the required
\(\forall(\mathcal D,h)\,\exists i\) interface.  No selector, fiber
probability, finite-support tape reduction, irrelevant catalog coordinate,
other policy, or invalid transcript enters the conclusion.
\end{proof}

\subsection{Finite-Distribution Minimax and Simultaneous Separation}
\label{app:step-003}

For a nonempty finite set \(F\), write
\[
\Delta_F:=\left\{p\in[0,1]^F:\sum_{x\in F}p_x=1\right\},
\qquad [L]:=\{1,\ldots,L\}.
\]
Let \(e_i\) denote the \(i\)-th vertex of \(\Delta_L\), and let \(d_x\)
denote the point-mass vertex of \(\Delta_F\).

\begin{lemma}[Finite-support instances lower-bound the distribution game]
\label{lem:step-003-game-lower-bound}
Under Assumptions~\ref{assump:sq-parameter-regime},
\ref{assump:universal-adversarial-sq}, and
\ref{assump:canonical-rounded-output-catalog}, together with
Proposition~\ref{prop:step-002-exact-correlation}, fix
\(h\in\mathcal H\) and a nonempty finite \(F\subseteq\mathcal X\).  Define
\begin{equation}
A\in\mathbb R^{F\times L},
\qquad
A_{xi}:=h(x)g_i(x).
\label{eq:step-003-payoff-matrix}
\end{equation}
Then every \(p\in\Delta_F\) satisfies
\[
\max_{i\in[L]}p^{T}Ae_i\ge\rho_{\varepsilon},
\]
and the minimum exists and obeys
\begin{equation}
\min_{p\in\Delta_F}\max_{i\in[L]}p^{T}Ae_i
\ge\rho_{\varepsilon}.
\label{eq:step-003-left-value}
\end{equation}
\end{lemma}

\begin{proof}
Fix \(p\in\Delta_F\).  Define the finitely supported probability law
\begin{equation}
\mathcal D_p:=\sum_{x\in F}p_x\delta_x,
\qquad
\mathbb E_{z\sim\mathcal D_p}[u(z)]
:=\sum_{x\in F}p_xu(x)
\label{eq:step-003-finite-law}
\end{equation}
for every pointwise function used below.  Coordinates with \(p_x=0\) cause
no difficulty; the support is merely contained in \(F\).  This finite-sum
definition introduces no singleton-measurability condition.

Proposition~\ref{prop:step-002-exact-correlation} applies because it ranges
over every distribution, not only full-support or nonatomic laws.  It gives
an index \(i(p)\in[L]\) such that
\begin{align}
p^{T}Ae_{i(p)}
&=\sum_{x\in F}p_xA_{x,i(p)} \notag\\
&=\sum_{x\in F}p_xh(x)g_{i(p)}(x) \notag\\
&=\mathbb E_{x\sim\mathcal D_p}[h(x)g_{i(p)}(x)] \notag\\
&\ge\rho_{\varepsilon}.
\label{eq:step-003-finite-correlation}
\end{align}
Every expectation here is the displayed finite sum, so this proof instance
adds no topology or measurability condition on \(F\), \(p\), or the catalog
coordinates.  It also does not narrow the universal premise to finite-support
distributions.

Since \(i(p)\) is among the finitely many maximized indices,
\begin{equation}
\max_{i\in[L]}p^{T}Ae_i
\ge p^{T}Ae_{i(p)}
\ge\rho_{\varepsilon}.
\label{eq:step-003-every-p}
\end{equation}
No common index for different \(p\)'s is asserted.  The map
\(p\mapsto\max_i p^{T}Ae_i\) is a finite maximum of continuous affine
functions.  The nonempty simplex \(\Delta_F\) is closed and bounded in a
finite-dimensional Euclidean space, hence compact, so this continuous map
has a minimizer \(p_*\).  Evaluating \eqref{eq:step-003-every-p} at \(p_*\)
gives
\[
\min_{p\in\Delta_F}\max_{i\in[L]}p^{T}Ae_i
=\max_{i\in[L]}p_*^{T}Ae_i
\ge\rho_{\varepsilon}.
\]
\end{proof}

\begin{lemma}[Catalog-simplex maximization occurs at a column vertex]
\label{lem:step-003-column-vertices}
Under Assumption~\ref{assump:canonical-rounded-output-catalog} and the matrix
construction \eqref{eq:step-003-payoff-matrix}, if \(F\ne\varnothing\) is
finite and \(p\in\Delta_F\), then
\begin{equation}
\max_{w\in\Delta_L}p^{T}Aw
=\max_{i\in[L]}p^{T}Ae_i,
\label{eq:step-003-column-identity}
\end{equation}
and both sides are attained.
\end{lemma}

\begin{proof}
Fix \(p\) and put \(c_i:=p^{T}Ae_i\).  For every \(w\in\Delta_L\),
\begin{align}
p^{T}Aw
&=\sum_{i=1}^Lw_i p^{T}Ae_i
=\sum_{i=1}^Lw_ic_i \notag\\
&\le\sum_{i=1}^Lw_i\max_{j\in[L]}c_j
=\max_{j\in[L]}c_j.
\label{eq:step-003-column-upper}
\end{align}
Because \([L]\) is finite and nonempty, choose
\(i_*\in\operatorname*{arg\,max}_i c_i\).  The vertex
\(e_{i_*}\in\Delta_L\) gives
\begin{equation}
p^{T}Ae_{i_*}=c_{i_*}=\max_{i\in[L]}c_i.
\label{eq:step-003-column-attainment}
\end{equation}
Equations \eqref{eq:step-003-column-upper} and
\eqref{eq:step-003-column-attainment} prove the identity and attainment.
Zero coordinates of \(w\) and duplicate columns are allowed.
\end{proof}

\begin{lemma}[Distribution-simplex minimization occurs at a point mass]
\label{lem:step-003-row-vertices}
Under the matrix construction \eqref{eq:step-003-payoff-matrix}, if
\(F\ne\varnothing\) is finite and \(w\in\Delta_L\), then
\begin{equation}
\min_{p\in\Delta_F}p^{T}Aw
=\min_{x\in F}(Aw)_x,
\label{eq:step-003-row-identity}
\end{equation}
and both sides are attained.
\end{lemma}

\begin{proof}
For every \(p\in\Delta_F\),
\begin{align}
p^{T}Aw
&=\sum_{x\in F}p_x(Aw)_x \notag\\
&\ge\sum_{x\in F}p_x\min_{y\in F}(Aw)_y
=\min_{y\in F}(Aw)_y.
\label{eq:step-003-row-lower}
\end{align}
Since \(F\) is finite and nonempty, choose
\(x_*\in\operatorname*{arg\,min}_{x\in F}(Aw)_x\).  Its point mass
\(d_{x_*}\in\Delta_F\) satisfies
\begin{equation}
d_{x_*}^{T}Aw=(Aw)_{x_*}=\min_{x\in F}(Aw)_x.
\label{eq:step-003-row-attainment}
\end{equation}
This proves the identity and attainment.  In particular, boundary points of
\(\Delta_F\), whose other coordinates vanish, are included.
\end{proof}

\begin{proposition}[Exact finite minimax separator]
\label{prop:step-003-finite-separator}
Under Assumptions~\ref{assump:sq-parameter-regime},
\ref{assump:universal-adversarial-sq}, and
\ref{assump:canonical-rounded-output-catalog}, together with
Proposition~\ref{prop:step-002-exact-correlation} and
Lemmas~\ref{lem:step-003-game-lower-bound},
\ref{lem:step-003-column-vertices}, and
\ref{lem:step-003-row-vertices}, fix \(h\in\mathcal H\) and a nonempty
finite \(F\subseteq\mathcal X\), and let
\(A_{xi}=h(x)g_i(x)\).  Then
\begin{equation}
\min_{p\in\Delta_F}\max_{i\in[L]}p^{T}Ae_i
=\max_{w\in\Delta_L}\min_{x\in F}(Aw)_x
\ge\rho_{\varepsilon}.
\label{eq:step-003-minimax-separator}
\end{equation}
The maximum on the right is attained by some
\(w_{h,F}\in\Delta_L\), and
\begin{equation}
\forall x\in F,
\qquad
h(x)\sum_{i=1}^Lw_{h,F,i}g_i(x)
\ge\rho_{\varepsilon}.
\label{eq:step-003-finite-margin}
\end{equation}
\end{proposition}

\begin{proof}
We first record the cited minimax statement in the orientation used here.
Sion's Theorem~3.4 \citep{Sion1958} says that if \(M\) is the maximizing
domain, \(N\) is the minimizing domain, both are nonempty compact convex
sets, and a real payoff \(f:M\times N\to\mathbb R\) is upper semicontinuous
and quasi-concave in its first variable and lower semicontinuous and
quasi-convex in its second, then
\begin{equation}
\sup_{w\in M}\inf_{p\in N}f(w,p)
=\inf_{p\in N}\sup_{w\in M}f(w,p).
\label{eq:step-003-sion-statement}
\end{equation}

Take \(M=\Delta_L\), \(N=\Delta_F\), and
\(f(w,p):=p^{T}Aw\).  The conditions \(L\ge1\) and
\(F\ne\varnothing\) make the simplices nonempty.  Each simplex is convex,
closed, and bounded in a finite-dimensional Euclidean space, hence compact.
For fixed \(p\), the function \(w\mapsto p^{T}Aw\) is real-valued,
continuous, and affine, hence quasi-concave.  For fixed \(w\), the function
\(p\mapsto p^{T}Aw\) is real-valued, continuous, and affine, hence
quasi-convex.  Thus the variables, topology, convexity, semicontinuity, and
quasi-concavity/quasi-convexity conventions all match
\eqref{eq:step-003-sion-statement}.  The theorem gives only
\begin{equation}
\sup_{w\in\Delta_L}\inf_{p\in\Delta_F}p^{T}Aw
=\inf_{p\in\Delta_F}\sup_{w\in\Delta_L}p^{T}Aw.
\label{eq:step-003-sion-applied}
\end{equation}

Attainment is checked separately.  By
Lemma~\ref{lem:step-003-row-vertices}, the left inner infimum is the function
\(w\mapsto\min_{x\in F}(Aw)_x\), a finite minimum of continuous affine
functions and therefore continuous.  Compactness of the nonempty
\(\Delta_L\) gives a maximizer \(w_{h,F}\).  By
Lemma~\ref{lem:step-003-column-vertices}, the right inner supremum is
\(p\mapsto\max_{i\in[L]}p^{T}Ae_i\), a finite maximum of continuous affine
functions.  Its minimum on compact \(\Delta_F\) is attained, as also shown
in Lemma~\ref{lem:step-003-game-lower-bound}.  Hence the suprema and infima
in \eqref{eq:step-003-sion-applied} may be written as maxima and minima.

Substituting the two independently proved vertex identities and then using
Lemma~\ref{lem:step-003-game-lower-bound} gives
\begin{align}
\max_{w\in\Delta_L}\min_{x\in F}(Aw)_x
&=\max_{w\in\Delta_L}\min_{p\in\Delta_F}p^{T}Aw \notag\\
&=\min_{p\in\Delta_F}\max_{w\in\Delta_L}p^{T}Aw \notag\\
&=\min_{p\in\Delta_F}\max_{i\in[L]}p^{T}Ae_i \notag\\
&\ge\rho_{\varepsilon}.
\label{eq:step-003-game-chain}
\end{align}
Reading the equality in the opposite direction gives exactly the orientation
in \eqref{eq:step-003-minimax-separator}: the distribution player is the
outer minimizer and the catalog-mixture player is the outer maximizer.
Sion's theorem supplies the middle equality only;
Lemma~\ref{lem:step-003-game-lower-bound} supplies positivity, and
Lemmas~\ref{lem:step-003-column-vertices} and
\ref{lem:step-003-row-vertices} supply the endpoint identities.

For the maximizer \(w_{h,F}\), \eqref{eq:step-003-game-chain} gives
\(\min_{x\in F}(Aw_{h,F})_x\ge\rho_{\varepsilon}\).  Therefore, for each
\(x\in F\),
\begin{equation}
(Aw_{h,F})_x
=\sum_{i=1}^LA_{xi}w_{h,F,i}
=h(x)\sum_{i=1}^Lw_{h,F,i}g_i(x)
\ge\rho_{\varepsilon}.
\label{eq:step-003-matrix-pointwise-identity}
\end{equation}
The matrix payoff and the exported pointwise score are identical; there is no
surrogate object or margin loss.

For a singleton \(F\), \(\Delta_F\) contains only its point mass.  For
\(L=1\), \(\Delta_L=\{e_1\}\).  Duplicate columns preserve every
calculation, and zero weights are permitted.  If \(\varepsilon=0\), then
\(\rho_{\varepsilon}=1\).  The game value is at least one, while every entry
of \(A\) belongs to \(\{-1,+1\}\), so every convex-combination payoff is at
most one.  Hence the game value and every finite-set margin are exactly one.
\end{proof}

\begin{lemma}[The empty finite restriction is vacuously feasible]
\label{lem:step-003-empty-restriction}
Under Assumptions~\ref{assump:sq-parameter-regime} and
\ref{assump:canonical-rounded-output-catalog}, if \(h\in\mathcal H\) and
\(F=\varnothing\), then \(e_1\in\Delta_L\) and
\[
\forall x\in F,
\qquad
h(x)\sum_{i=1}^L(e_1)_ig_i(x)\ge\rho_{\varepsilon}
\]
holds vacuously.
\end{lemma}

\begin{proof}
Assumption~\ref{assump:canonical-rounded-output-catalog} gives \(L\ge1\), so
the first simplex vertex \(e_1\) exists and belongs to \(\Delta_L\).  There
is no \(x\in\varnothing\) at which the inequality can fail.  This branch
defines no probability simplex or probability distribution indexed by the
empty set, takes no minimum over an empty index set, and does not invoke
minimax.  Together with Proposition~\ref{prop:step-003-finite-separator}, it
proves exact \(\rho_{\varepsilon}\)-feasibility on every finite restriction,
without asserting compatible choices for different finite sets.
\end{proof}

\begin{proof}[Conclusion of the finite-separation argument]
Fix \(h\in\mathcal H\) and finite \(F\subseteq\mathcal X\).  If
\(F=\varnothing\), Lemma~\ref{lem:step-003-empty-restriction} supplies a
weight in \(\Delta_L\) and the constraints are vacuous.  If
\(F\ne\varnothing\), Lemma~\ref{lem:step-003-game-lower-bound} gives the
positive distribution-game value, while Lemmas
\ref{lem:step-003-column-vertices} and
\ref{lem:step-003-row-vertices} prove the two simplex identities.
Proposition~\ref{prop:step-003-finite-separator} combines these facts with
the fully checked Sion equality and independent attainment to obtain
\[
\min_{p\in\Delta_F}\max_{i\in[L]}p^{T}Ae_i
=\max_{w\in\Delta_L}\min_{x\in F}(Aw)_x
\ge\rho_{\varepsilon},
\qquad A_{xi}=h(x)g_i(x),
\]
and a maximizer \(w_{h,F}\) satisfying
\[
h(x)\sum_{i=1}^Lw_{h,F,i}g_i(x)
\ge\rho_{\varepsilon}
\qquad(x\in F).
\]
Since \(h\) and \(F\) were arbitrary, both branches give exact finite
feasibility.  No compatibility between weights for different \(F\)'s is
asserted or needed.
\end{proof}

\subsection{Fixed-Simplex Globalization}
\label{app:step-004}

\begin{lemma}[The fixed catalog simplex is nonempty and compact]
\label{lem:step-004-simplex-compact}
Under Assumption~\ref{assump:canonical-rounded-output-catalog}, the simplex
\[
\Delta_L=\left\{w\in[0,1]^L:\sum_{i=1}^Lw_i=1\right\}
\subseteq\mathbb R^L
\]
is nonempty and compact.
\end{lemma}

\begin{proof}
Assumption~\ref{assump:canonical-rounded-output-catalog} gives a finite
integer \(L\ge1\).  Hence \(e_1=(1,0,\ldots,0)\) belongs to \(\Delta_L\),
proving nonemptiness.

For each coordinate, the condition \(0\le w_i\le1\) defines a closed subset
of \(\mathbb R^L\).  The map \(w\mapsto\sum_{i=1}^Lw_i\) is continuous, so
the affine hyperplane on which the sum equals one is closed.  Their finite
intersection is \(\Delta_L\), which is therefore closed.  Moreover,
\begin{equation}
\lVert w\rVert_2^2
=\sum_{i=1}^Lw_i^2
\le\sum_{i=1}^Lw_i
=1
\qquad (w\in\Delta_L),
\label{eq:step-004-simplex-bound}
\end{equation}
where \(w_i^2\le w_i\) follows from \(0\le w_i\le1\).  Thus the simplex is
bounded.  Since \(L\) is finite, the Heine--Borel theorem gives compactness.
For \(L=1\), this reads \(\Delta_1=\{(1)\}\).
\end{proof}

\begin{lemma}[Each pointwise margin constraint is relatively closed]
\label{lem:step-004-constraint-closed}
Under Assumptions~\ref{assump:sq-parameter-regime} and
\ref{assump:canonical-rounded-output-catalog}, fix \(h\in\mathcal H\) and
\(x\in\mathcal X\), and define
\begin{equation}
\ell_{h,x}(w):=h(x)\sum_{i=1}^Lw_i g_i(x),
\qquad
C_{h,x}:=\{w\in\Delta_L:\ell_{h,x}(w)\ge\rho_{\varepsilon}\}.
\label{eq:step-004-point-constraint}
\end{equation}
Then \(\ell_{h,x}:\mathbb R^L\to\mathbb R\) is a continuous linear
functional and \(C_{h,x}\) is relatively closed in \(\Delta_L\).
\end{lemma}

\begin{proof}
The coefficients \(h(x)g_i(x)\) lie in \(\{-1,+1\}\).  For
\(u,v\in\mathbb R^L\),
\begin{align}
|\ell_{h,x}(u)-\ell_{h,x}(v)|
&=\left|\sum_{i=1}^Lh(x)g_i(x)(u_i-v_i)\right| \notag\\
&\le\sum_{i=1}^L|u_i-v_i|
\le L\lVert u-v\rVert_2.
\label{eq:step-004-functional-lipschitz}
\end{align}
Thus \(\ell_{h,x}\) is Lipschitz and continuous.  The ray
\([\rho_{\varepsilon},\infty)\) is closed in \(\mathbb R\), and
\begin{equation}
C_{h,x}
=\Delta_L\cap\ell_{h,x}^{-1}([\rho_{\varepsilon},\infty))
=\bigl(\ell_{h,x}|_{\Delta_L}\bigr)^{-1}
([\rho_{\varepsilon},\infty)).
\label{eq:step-004-closed-inverse-image}
\end{equation}
Hence \(C_{h,x}\) is closed in the subspace topology of \(\Delta_L\).  The
proof uses \(x\) only as an index and assumes no topology or sigma-algebra on
\(\mathcal X\).  Duplicate catalog functions do not affect the calculation.
\end{proof}

\begin{lemma}[Finite separators give the finite-intersection property]
\label{lem:step-004-finite-intersections}
Under Assumptions~\ref{assump:sq-parameter-regime},
\ref{assump:universal-adversarial-sq}, and
\ref{assump:canonical-rounded-output-catalog}, together with
Proposition~\ref{prop:step-003-finite-separator},
Lemma~\ref{lem:step-003-empty-restriction}, and
Lemma~\ref{lem:step-004-simplex-compact}, fix \(h\in\mathcal H\) and define
\(C_{h,x}\) by \eqref{eq:step-004-point-constraint}.  Then every finite
\(F\subseteq\mathcal X\), including \(F=\varnothing\), satisfies
\begin{equation}
\bigcap_{x\in F}C_{h,x}\ne\varnothing,
\label{eq:step-004-fip}
\end{equation}
where the empty-subfamily intersection is \(\Delta_L\).
\end{lemma}

\begin{proof}
If \(F\ne\varnothing\),
Proposition~\ref{prop:step-003-finite-separator} gives
\(w_{h,F}\in\Delta_L\) such that
\[
h(x)\sum_{i=1}^Lw_{h,F,i}g_i(x)
\ge\rho_{\varepsilon}
\qquad (x\in F).
\]
By \eqref{eq:step-004-point-constraint}, this is precisely
\[
w_{h,F}\in\bigcap_{x\in F}C_{h,x}.
\]

If \(F=\varnothing\), the intersection of the empty subfamily of subsets of
\(\Delta_L\) is, by definition, \(\Delta_L\).  It contains \(e_1\) by
Lemma~\ref{lem:step-003-empty-restriction}, while
Lemma~\ref{lem:step-004-simplex-compact} independently proves nonemptiness.
A singleton subfamily is covered by the nonempty branch.  The witness may
depend arbitrarily on \(F\); no compatible family
\(F\mapsto w_{h,F}\) is selected or assumed.
\end{proof}

\begin{lemma}[Closed-set finite-intersection principle]
\label{lem:step-004-compact-fip}
Let \(Y\) be a compact topological space, let \(J\) be an arbitrary index
set, and let \(\{D_j:j\in J\}\) be a family of subsets closed in \(Y\).
Adopt the convention \(\bigcap_{j\in\varnothing}D_j=Y\).  If
\[
\bigcap_{j\in E}D_j\ne\varnothing
\qquad\text{for every finite }E\subseteq J,
\]
including \(E=\varnothing\), then
\[
\bigcap_{j\in J}D_j\ne\varnothing.
\]
\end{lemma}

\begin{proof}
The hypothesis for \(E=\varnothing\) says \(Y\ne\varnothing\).  If
\(J=\varnothing\), the total intersection is \(Y\), so the conclusion is
immediate.  Otherwise, suppose for contradiction that
\begin{equation}
\bigcap_{j\in J}D_j=\varnothing.
\label{eq:step-004-empty-total-intersection}
\end{equation}
For \(j\in J\), let \(U_j:=Y\setminus D_j\).  Each \(U_j\) is open in
\(Y\).  Taking complements in
\eqref{eq:step-004-empty-total-intersection} gives
\[
Y=\bigcup_{j\in J}U_j,
\]
so the \(U_j\)'s form an open cover.  Compactness yields a finite
\(E\subseteq J\) such that
\[
Y=\bigcup_{j\in E}U_j.
\]
Taking complements inside \(Y\) and applying De Morgan's law gives
\begin{equation}
\bigcap_{j\in E}D_j
=Y\setminus\bigcup_{j\in E}U_j
=\varnothing,
\label{eq:step-004-finite-contradiction}
\end{equation}
contradicting the assumed nonemptiness of every finite intersection.  The
argument applies to every cardinality of \(J\) and uses neither a sequence
nor sequential compactness.
\end{proof}

\begin{proposition}[One fixed-simplex separator works on the arbitrary domain]
\label{prop:step-004-global-separator}
Under Assumptions~\ref{assump:sq-parameter-regime},
\ref{assump:universal-adversarial-sq}, and
\ref{assump:canonical-rounded-output-catalog}, together with
Proposition~\ref{prop:step-003-finite-separator},
Lemma~\ref{lem:step-003-empty-restriction}, and
Lemmas~\ref{lem:step-004-simplex-compact},
\ref{lem:step-004-constraint-closed},
\ref{lem:step-004-finite-intersections}, and
\ref{lem:step-004-compact-fip}, for every \(h\in\mathcal H\) there is a
\(w_h\in\Delta_L\) such that
\begin{equation}
\forall x\in\mathcal X,
\qquad
h(x)\sum_{i=1}^Lw_{h,i}g_i(x)
\ge\rho_{\varepsilon}=1-2\varepsilon.
\label{eq:step-004-global-margin}
\end{equation}
\end{proposition}

\begin{proof}
Fix \(h\in\mathcal H\).  By
Lemma~\ref{lem:step-004-simplex-compact}, the space
\(Y:=\Delta_L\) is nonempty and compact.  By
Lemma~\ref{lem:step-004-constraint-closed}, every
\(D_x:=C_{h,x}\), \(x\in\mathcal X\), is closed in this same space.  By
Lemma~\ref{lem:step-004-finite-intersections},
\[
\bigcap_{x\in F}D_x\ne\varnothing
\qquad\text{for every finite }F\subseteq\mathcal X,
\]
including the empty subfamily.  Lemma~\ref{lem:step-004-compact-fip}
therefore applies with \(J=\mathcal X\), and gives
\begin{equation}
\bigcap_{x\in\mathcal X}C_{h,x}\ne\varnothing.
\label{eq:step-004-total-intersection}
\end{equation}
Choose \(w_h\) from this intersection.  For each \(x\), membership in
\(C_{h,x}\) unfolds exactly to \eqref{eq:step-004-global-margin}.  The finite
and global inequalities use the identical catalog, target, simplex, point
evaluation, and threshold; the transfer has zero residual.

If \(\mathcal X=\varnothing\), the total intersection is the empty-family
intersection \(\Delta_L\), and \(e_1\) is an explicit choice; the pointwise
conclusion is vacuous.  If \(\mathcal X\) is finite, countably infinite, or
uncountable, the same arbitrary-index-set lemma applies without a topology,
ordering, sigma-algebra, or countability condition on \(\mathcal X\).  If
\(L=1\), the simplex is \(\{(1)\}\), and finite feasibility puts this sole
weight in every finite intersection.  Duplicate catalog functions alter
neither the ambient simplex, closedness, nor finite feasibility.

If \(\varepsilon=0\), every constraint has threshold one.  For any
\(w\in\Delta_L\), the value
\(h(x)\sum_iw_i g_i(x)\) is a convex combination of numbers in
\(\{-1,+1\}\), so it is at most one.  Thus every nonvacuous global margin is
exactly one, not one minus a compactness residual.  At no point is a
compatible family of finite witnesses selected: all constraints live in the
single fixed simplex, and no varying simplex, limiting simplex, subsequence,
net, or topology on the index set is introduced.
\end{proof}

\begin{proof}[Conclusion of the fixed-simplex globalization argument]
For an arbitrary \(h\), Lemma~\ref{lem:step-004-simplex-compact} supplies the
one nonempty compact ambient simplex,
Lemma~\ref{lem:step-004-constraint-closed} supplies the exact closed
pointwise constraints, and Lemma~\ref{lem:step-004-finite-intersections}
translates Propositions~\ref{prop:step-003-finite-separator} and
Lemma~\ref{lem:step-003-empty-restriction} into their finite-intersection
property.  Lemma~\ref{lem:step-004-compact-fip} then yields the total
intersection, and Proposition~\ref{prop:step-004-global-separator} gives one
\(w_h\) satisfying \eqref{eq:step-004-global-margin} on the entire arbitrary
domain.  Since \(h\) was arbitrary, the quantifiers are
\(\forall h\,\exists w_h\,\forall x\), and the threshold remains
\(\rho_{\varepsilon}\) with zero loss.
\end{proof}

\subsection{Exact Coordinates, Signs, and Dimension Closure}
\label{app:step-005}

\begin{lemma}[Exact pre-instance catalog coordinates]
\label{lem:step-005-coordinate-map}
Under Assumption~\ref{assump:canonical-rounded-output-catalog} and
Proposition~\ref{prop:step-001-canonical-membership}, let
\(\mathcal C_A^{\rho}=(g_1,\ldots,g_L)\) be the exact ordered catalog and
define
\begin{equation}
\phi:\mathcal X\to\mathbb R^L,
\qquad
\phi(x):=(g_1(x),\ldots,g_L(x)).
\label{eq:step-005-coordinate-map}
\end{equation}
Then \(\phi\) is fixed before and independently of the current distribution,
target, oracle policy, and learner tape.  For every
\(w=(w_1,\ldots,w_L)\in\mathbb R^L\) and \(x\in\mathcal X\),
\begin{equation}
\langle w,\phi(x)\rangle
=\sum_{i=1}^Lw_i g_i(x).
\label{eq:step-005-coordinate-identity}
\end{equation}
The identity is exact for \(L=1\), repeated catalog functions, and arbitrary
finite or infinite \(\mathcal X\).
\end{lemma}

\begin{proof}
Assumption~\ref{assump:canonical-rounded-output-catalog} fixes the ordered
list \((g_1,\ldots,g_L)\) after the protocol and \(\rho\) are fixed but
before the current distribution, target, oracle interaction, or tape.
Proposition~\ref{prop:step-001-canonical-membership} uses this same ordered
catalog as the exact function-valued interface on actual canonical
executions.  Thus \eqref{eq:step-005-coordinate-map} uses only pre-instance
catalog data and makes no selection depending on \(\mathcal D\), \(h\), a
policy, a reply, or a tape.  The map is deterministic, rather than sampled
from a random feature-map law.

For each \(x\), the binary catalog gives
\(\phi(x)\in\{-1,+1\}^L\subseteq\mathbb R^L\).  By the Euclidean inner
product and the coordinate definition,
\[
\langle w,\phi(x)\rangle
=\sum_{i=1}^Lw_i\phi_i(x)
=\sum_{i=1}^Lw_i g_i(x),
\]
with zero residual.  Repeated functions simply give repeated legal
coordinates and require no linear independence.  For \(L=1\), the formula
is \(\langle(w_1),(g_1(x))\rangle=w_1g_1(x)\).  The calculation is pointwise
and finite-dimensional, so the cardinality of \(\mathcal X\) is irrelevant.
\end{proof}

\begin{proposition}[Exact target signs from the global simplex separator]
\label{prop:step-005-exact-signs}
Under Assumptions~\ref{assump:sq-parameter-regime},
\ref{assump:universal-adversarial-sq}, and
\ref{assump:canonical-rounded-output-catalog}, together with
Proposition~\ref{prop:step-004-global-separator} and
Lemma~\ref{lem:step-005-coordinate-map}, the one fixed map \(\phi\) satisfies
\begin{equation}
\forall h\in\mathcal H\ \exists w_h\in\Delta_L\ \forall x\in\mathcal X,
\qquad
h(x)\langle w_h,\phi(x)\rangle
\ge1-2\varepsilon>\frac12>0.
\label{eq:step-005-exact-signs}
\end{equation}
Only \(h\) may affect the weight.  The assertion includes empty
\(\mathcal H\), empty \(\mathcal X\), \(L=1\), repeated coordinates, and
infinite domains.
\end{proposition}

\begin{proof}
Assumption~\ref{assump:sq-parameter-regime} gives
\(0\le\varepsilon<1/4\), and therefore
\begin{equation}
2\varepsilon<\frac12
\quad\Longrightarrow\quad
\rho_{\varepsilon}=1-2\varepsilon>\frac12>0.
\label{eq:step-005-positive-margin}
\end{equation}
Fix \(h\in\mathcal H\).  Proposition
\ref{prop:step-004-global-separator} gives \(w_h\in\Delta_L\) such that
simultaneously for every \(x\in\mathcal X\),
\begin{equation}
h(x)\sum_{i=1}^Lw_{h,i}g_i(x)
\ge\rho_{\varepsilon}=1-2\varepsilon.
\label{eq:step-005-mixture-margin}
\end{equation}
Using \eqref{eq:step-005-coordinate-identity} with \(w=w_h\) yields
\begin{align}
h(x)\langle w_h,\phi(x)\rangle
&=h(x)\sum_{i=1}^Lw_{h,i}g_i(x) \notag\\
&\ge1-2\varepsilon
>\frac12>0.
\label{eq:step-005-score-margin}
\end{align}
The equality is coordinate-for-coordinate, and the inequality is exactly
\eqref{eq:step-005-mixture-margin}; no approximate sign or distributional
exceptional set appears.

The catalog and \(\phi\) were fixed before the target was chosen.  The
proposition producing \(w_h\) fixes only \(h\) and uses the already fixed
catalog; its conclusion contains no current distribution, policy, reply, or
tape.  The quantifier order is therefore one common \(\phi\), followed by a
target-specific \(w_h\), followed by all points.  No measurable selector
\(h\mapsto w_h\) is required.

If \(\mathcal H=\varnothing\), the target quantifier is vacuous.  If
\(\mathcal X=\varnothing\) and \(h\in\mathcal H\), Proposition
\ref{prop:step-004-global-separator} supplies a point of nonempty
\(\Delta_L\), and the pointwise clause is vacuous.  If \(L=1\), membership
in \(\Delta_1\) forces \(w_h=(1)\).  Duplicate coordinates do not change
the exact identity, and infinite domains are covered pointwise by the global
separator.
\end{proof}

\begin{proposition}[Deterministic dimension closure with the catalog bound]
\label{prop:step-005-dimension-closure}
Under Assumptions~\ref{assump:sq-parameter-regime},
\ref{assump:universal-adversarial-sq}, and
\ref{assump:canonical-rounded-output-catalog}, together with
Proposition~\ref{prop:step-005-exact-signs} and the definition
\eqref{eq:dc-definition},
\begin{equation}
\operatorname{dc}(\mathcal H)
\le L
\le B\left(1+\frac{m}{\tau^2}\right)^k.
\label{eq:step-005-dimension-closure}
\end{equation}
If \(\mathcal H=\varnothing\) or \(\mathcal X=\varnothing\), then
\(\operatorname{dc}(\mathcal H)=0\).  The specializations \(m=0\),
\(\varepsilon=0\), \(L=1\), \(B=1\), \(k=1\), every \(\tau>0\),
infinite domains, duplicate coordinates, and arbitrary including nonatomic
tape spaces retain the literal conclusions.
\end{proposition}

\begin{proof}
Suppose first that \(\mathcal H\ne\varnothing\) and
\(\mathcal X\ne\varnothing\).  Proposition
\ref{prop:step-005-exact-signs} supplies one map
\(\phi:\mathcal X\to\mathbb R^L\) such that for every
\(h\in\mathcal H\) there is \(w_h\in\mathbb R^L\) with
\[
\forall x\in\mathcal X,
\qquad
h(x)\langle w_h,\phi(x)\rangle>0.
\]
Thus \(d=L\), \(\psi=\phi\), and \(u_h=w_h\) are admissible in
\eqref{eq:dc-definition}.  Since \(\operatorname{dc}(\mathcal H)\) is the
least admissible dimension,
\begin{equation}
\operatorname{dc}(\mathcal H)\le L.
\label{eq:step-005-dc-at-most-l}
\end{equation}

Now suppose \(\mathcal H=\varnothing\) or \(\mathcal X=\varnothing\).  Let
\(d=0\) and let \(\psi_0:\mathcal X\to\mathbb R^0\) be the unique map.  If
\(\mathcal H=\varnothing\), the target quantifier in
\eqref{eq:dc-definition} is empty.  If \(\mathcal X=\varnothing\), choose
the unique \(u_h\in\mathbb R^0\) for each target; the point quantifier is
empty.  In either case the strict-sign condition is vacuous, so \(d=0\) is
admissible.  Since dimensions lie in \(\mathbb N_0\),
\begin{equation}
\operatorname{dc}(\mathcal H)=0.
\label{eq:step-005-empty-dimension}
\end{equation}
As \(L\ge1\), this also implies \eqref{eq:step-005-dc-at-most-l} in the
degenerate branches.

Assumption~\ref{assump:canonical-rounded-output-catalog} states literally
that
\begin{equation}
L\le B\left(1+\frac{m}{\tau^2}\right)^k.
\label{eq:step-005-literal-catalog-bound}
\end{equation}
Appending \eqref{eq:step-005-literal-catalog-bound} to
\eqref{eq:step-005-dc-at-most-l} proves
\eqref{eq:step-005-dimension-closure}.  No term is dropped, dominated, or
absorbed, and no hidden factor appears.

The boundary specializations are exact.  If \(m=0\), then
\[
L\le B\left(1+\frac{0}{\tau^2}\right)^k=B.
\]
If \(\varepsilon=0\), the lower margin in
\eqref{eq:step-005-score-margin} is one.  Since \(w_h\in\Delta_L\) and each
\(h(x)g_i(x)\in\{-1,+1\}\),
\[
h(x)\langle w_h,\phi(x)\rangle
\le\sum_{i=1}^Lw_{h,i}=1,
\]
so every nonvacuous margin is exactly one.  If \(L=1\), then
\(\Delta_1=\{(1)\}\) and \(\phi(x)=(g_1(x))\).  If \(B=1\), the catalog
bound is \(L\le(1+m/\tau^2)^k\); if \(k=1\), it is
\(L\le B(1+m/\tau^2)\).  If \(B=k=1\) and \(m=0\), then
\(1\le L\le1\), so \(L=1\).  Every \(\tau>0\) is allowed because
\(\tau^2>0\), and this argument imposes no further restriction.

The cardinality of \(\mathcal X\) does not enter the coordinate identity,
the pointwise margin, or the dimension implication.  Repeated functions are
legal coordinates, and no rank or distinctness premise is used.  The
conclusion is deterministic: the tape and its law do not occur in \(\phi\),
\(w_h\), or the final inequalities, so arbitrary and nonatomic tape spaces
require no almost-sure, in-probability, or expectation qualification.
\end{proof}

\begin{proof}[Conclusion of the exact-coordinate and dimension argument]
Proposition~\ref{prop:step-001-canonical-membership} and
Assumption~\ref{assump:canonical-rounded-output-catalog} provide the exact
ordered pre-instance catalog.  Lemma~\ref{lem:step-005-coordinate-map} uses
those functions to define the common map and proves
\(\langle w_h,\phi(x)\rangle=\sum_iw_{h,i}g_i(x)\).  For every target,
Proposition~\ref{prop:step-004-global-separator} supplies the same-catalog
mixture at margin \(\rho_{\varepsilon}\) on every point.  Proposition
\ref{prop:step-005-exact-signs} combines the two exact interfaces with
\(\varepsilon<1/4\) to yield
\[
\forall h\in\mathcal H\ \exists w_h\in\Delta_L\ \forall x\in\mathcal X,
\qquad
h(x)\langle w_h,\phi(x)\rangle
\ge1-2\varepsilon>\frac12>0.
\]
Proposition~\ref{prop:step-005-dimension-closure} applies the exact dimension
definition and the literal catalog-size inequality to obtain
\[
\operatorname{dc}(\mathcal H)
\le L
\le B\left(1+\frac{m}{\tau^2}\right)^k.
\]
The map is common and pre-instance, only the weight is target-specific, and
the empty, zero-budget, zero-error, singleton-catalog, unit-constant,
arbitrary-tolerance, infinite-domain, duplicate-coordinate, and arbitrary
tape-space cases remain covered exactly.
\end{proof}

\subsection{Proof of the Main Theorem}
\label{app:proof-main-theorem}

\begin{proof}[Proof of Theorem~\ref{thm:main}]
Lemma~\ref{lem:step-001-canonical-validity} and
Proposition~\ref{prop:step-001-canonical-membership} establish the legal
canonical execution and its exact catalog membership.  Lemma
\ref{lem:step-002-attained-low-risk} and
Proposition~\ref{prop:step-002-exact-correlation} turn the universal
expected-risk guarantee into an exact per-distribution catalog-correlation
witness.  Proposition~\ref{prop:step-003-finite-separator}, together with
Lemma~\ref{lem:step-003-empty-restriction}, gives every finite simultaneous
separator.  Proposition~\ref{prop:step-004-global-separator} globalizes
those constraints in the fixed simplex.  Finally, Lemma
\ref{lem:step-005-coordinate-map} and Propositions
\ref{prop:step-005-exact-signs} and
\ref{prop:step-005-dimension-closure} give exactly
\eqref{eq:main-margin}, \eqref{eq:main-dimension-bound}, and the stated
zero-dimensional empty cases.  These are the conclusions of
Theorem~\ref{thm:main} under its three numbered assumptions.
\end{proof}
